% part: intuitionistic-logic
% chapter: propositions-as-types
% section: introduction

\documentclass[../../../include/open-logic-section]{subfiles}

\begin{document}

\olfileid{int}{pty}{int}

\olsection{परिचय}

लॅम्डा कलन हे संगणनाचे एक साधे प्रतिमान आहे. ते चरांपासून
बनलेल्या पदांवर कार्य करते. $N$ आणि $M$ ही पदे असतील, तर
$(NM)$ हेही एक पद आहे (``$N$ चे~$M$ वर उपयोजन''). $N$ हे पद
असेल, तर $\lambd[x][N]$ हेही पद आहे. $N$ मध्ये~$x$ हे चर
असेल, तर $\lambd[x][N]$ मधील~$x$ ची संबंधित स्थाने
$\lambd[x]$ या संकारकाने बद्ध होतात; त्यामुळे ती मुक्त राहत
नाहीत. $(\lambd[x][N]M)$ या रूपातील पदाचे
$\Subst{N}{M}{x}$ मध्ये $\beta$-संकुचन होते असे म्हणतात
(म्हणजे~$N$ मधील~$x$ ची प्रत्येक मुक्त आस्थिती~$M$ ने
बदलल्यावर मिळणारे पद). $N$ मधील एखाद्या उपपदाचे
$\beta$-संकुचन करून~$N'$ मिळत असेल, तर~$N$ चे~$N'$ मध्ये
$\beta$-परिवर्तन होते, आणि आपण $N \redone N'$ असे लिहितो.
$\lambd$-कलनातील संगणन म्हणजे
$N \redone N_1 \redone N_2 \redone \dots$ असा क्रम.
हा क्रम~$N_k$ या अशा पदावर संपला की त्याच्या कोणत्याही
उपपदाचे $\beta$-संकुचन होऊ शकत नाही, तर त्या पदाला
\emph{सामान्य} किंवा~$N$ चे सामान्य रूप म्हणतात.
$\lambd$-कलनातील संगणनाचा विचार अशा क्रमांच्या रूपात करा.
सामान्य रूप मिळाल्यास संगणन थांबते, आणि ते सामान्य रूपच
संगणनाचे फल असते. हे साधे संगणन-प्रतिमान ट्यूरिंग-संपूर्ण
ठरते: प्रत्येक आंशिक पुनरावर्ती फलनाची गणना
$\lambd$-कलनातील एका पदाने करता येते.

हास्केल करी आणि विल्यम हॉवर्ड यांनी स्वतंत्रपणे नैसर्गिक
निगमन आणि $\lambd$-कलन यांतील जवळचे साम्य शोधले.
सहज अर्थाने $\lambd[x][N]$ हे पद एक फलन आहे: $x$ हा त्याचा
आर्ग्युमेंट, आणि~$x$ दिल्यावर मूल्य कसे मोजायचे हे~$N$
सांगते. म्हणून $(\lambd[x][N]M)$ म्हणजे
$\lambd[x][N]$ हे फलन~$M$ या आर्ग्युमेंटवर लावणे;
$\beta$-संकुचनाने त्याचे मूल्यांकन कसे करायचे ते कळते:
$N$ मध्ये~$x$ ऐवजी~$M$ ठेवायचे (आणि मिळालेल्या पदाचे
मूल्यांकन पुढे चालू ठेवायचे).
आता या फलनांना निश्चित प्रांत आणि मूल्यसंच आहेत असे समजा.
$\lambd[x][N]$ चा प्रांत~$A$ आणि मूल्यसंच~$B$ असेल, तर
$x$ ला फक्त~$A$ मधील आर्ग्युमेंट घेता येतील आणि~$N$ ने
$B$ मधील मूल्ये द्यावीत. त्यामुळे $\lambd[x][N]$ हे
$A \to B$ असे फलन आणि $M \in A$ असेल, तेव्हाच
$(\lambd[x][N]M)$ ला अर्थ असावा. ही माहिती
$\lambd$-कलनात घातली की \emph{प्रकारयुक्त} $\lambd$-कलन
मिळते. प्रकारयुक्त $\lambd$-कलनात प्रत्येक
$(\lambd[x][N]M)$ अभिव्यक्ती ग्राह्य नसते;
$\lambd[x][N]$ आणि~$M$ यांचे ``प्रकार'' जुळतात,
म्हणजे $\lambd[x][N]$ चा प्रकार $A \lif B$ आणि~$M$ चा
प्रकार~$A$ असतो, तेव्हाच ती ग्राह्य असते.
$x$ चा प्रकार~$A$ आणि~$N$ चा प्रकार~$B$ असेल,
तेव्हाच $\lambd[x][N]$ चा प्रकार $A \to B$ असतो.
सर्वसाधारणपणे प्रत्येक $(M_1M_2)$ अभिव्यक्ती ग्राह्य नसते.
$(M_1M_2)$ हे पद ग्राह्य होण्यासाठी~$M_1$ चा प्रकार $A \to B$
आणि~$M_2$ चा प्रकार~$A$ असावा;
अशा वेळी $(M_1M_2)$ चा प्रकार~$B$ असतो.

हे प्रकार-नियम आता नैसर्गिक निगमनातील
!!{derivation}s शी स्पष्टपणे जुळतात: प्रकारांना
!!{formula}s, आणि !!{derivation}s ना $\lambd$-पदे
अनुरूप असतात. उदाहरणार्थ, $!A \lif !B$ ची
!!a{derivation} $\lambd[\typeof{x}{!A}][N]$
या पदाशी जुळते; त्याचा फलन-प्रकार $!A \lif !B$ आहे.
नैसर्गिक निगमनाचे अनुमान-नियम प्रकारयुक्त लॅम्डा कलनातील
प्रकार-नियमांशी जुळतात; उदाहरणार्थ,
\begin{prooftree}
  \AxiomC{$\Discharge{!A}{x}$}
  \DeduceC{$!B$}
  \DischargeRule{\Intro\lif}{x}
  \UnaryInfC{$!A \lif !B$}
  \DisplayProof\bottomAlignProof
  \quad\text{याला अनुरूप आहे}\quad
  \Axiom$x: !A \fCenter N: !B$
  \UnaryInf$ \fCenter \lambd[\typeof{x}{!A}][N] : !A \lif !B$
\end{prooftree}
उजवीकडील नियमाचा अर्थ असा की~$x$ चा प्रकार~$!A$ आणि
$N$ चा प्रकार~$!B$ असेल, तर
$\lambd[\typeof{x}{!A}][N]$ चा प्रकार~$!A \lif !B$ असतो.

$\Elim\lif$ हा नियम संयुक्त पदांसाठीच्या प्रकार-नियमाशी
जुळतो; म्हणजे,
\[
  \AxiomC{$!A \lif !B$}
  \AxiomC{$!A$}
  \RightLabel{\Elim\lif}
  \BinaryInfC{$!B$}
  \DisplayProof
  \quad\text{याला अनुरूप आहे}\quad
  \Axiom$\fCenter M_1: !A \lif !B$
  \Axiom$\fCenter M_2: !A$
  \BinaryInf$ \fCenter (M_1M_2) : !B$
  \DisplayProof
\]
\Intro\lif{} या नियमानंतर लगेच \Elim\lif{} नियम
आला, तर !!{derivation} सोपी करता येते:
\begin{gather*}
  \AxiomC{$\Discharge{!A}{x}$}
  \DeduceC{$!B$}
  \DischargeRule{\Intro\lif}{x}
  \UnaryInfC{$!A \lif !B$}
  \AxiomC{}
  \DeduceC{$!A$}
  \RightLabel{\Elim\lif}
  \BinaryInfC{$!B$}
  \DisplayProof
  \quad\redone\quad
  \AxiomC{}
  \DeduceC{$!A$}
  \DeduceC{$!B$}
  \DisplayProof
\end{gather*}
हे लॅम्डा पदांच्या $\beta$-संकुचनाशी जुळते:
\[
(\lambd[\typeof{x}{!A}][N]M) \quad\redone\quad
\Subst{N}{M}{x}.
\]

अशीच अनुरूपता $\land$ च्या नियमांचा ``गुणाकार''
प्रकारांशी आणि $\lor$ च्या नियमांचा ``बेरीज'' प्रकारांशीही आहे.

साध्या प्रकारयुक्त लॅम्डा कलनातील पदे आणि नैसर्गिक निगमनातील
!!{derivation}s यांतील या अनुरूपतेला ``करी--हॉवर्ड'',
``सिद्धता म्हणजे कार्यक्रम'', किंवा ``विधाने म्हणजे प्रकार''
अनुरूपता म्हणतात. !!{formula}s (विधाने) प्रकारांशी आणि
सिद्धता पदांशी जुळतात; त्याशिवाय इतर अनुरूपता पुढीलप्रमाणे:
\begin{center}
  \begin{tabular}{c c c}
    तर्कशास्त्र & कार्यक्रम \\
    \hline
    विधान/!!{formula} & प्रकार \\
    सिद्धता & पद \\
    गृहीतक & चर \\
    निरसित गृहीतक & बद्ध चर \\
    अनिरसित गृहीतक & मुक्त चर \\
    $!A \lif !B$ & फलन-प्रकार \\
    $!A \land !B$ & गुणाकार-प्रकार \\
    $!A \lor !B$ & बेरीज-प्रकार \\
    $\lfalse$ & रिक्त प्रकार \\
    \hline
  \end{tabular}
\end{center}

करी--हॉवर्ड अनुरूपता ही स्वयंचलित सिद्धता-साहाय्यकांची आणि
कार्यक्रमांसाठीच्या प्रकार-तपासकांची एक आधारभूत कल्पना आहे.
कारण एखादे विधान सिद्ध करणारी सिद्धता तपासणे (जसे आपण वर केले)
म्हणजे एखाद्या कार्यक्रमाला (पदाला) जाहीर केलेला प्रकार आहे का
हे तपासणे होय.
\end{document}
